\subsection{Random products}\label{subsec:random}

For the smaller moduli, two independent products of sampled primes supply
the cancellation.  Their integer values are smaller than a suitable
divisor of the modulus, so unique factorization controls their residue
distributions.

\begin{lemma}\label{lem:random}
Let \(K=100\), \(D_0=100000\), and \(m=\lfloor S/\log S\rfloor\).
Let \(\mathcal P\) be a set of primes in \([S^K,2S^K]\) with
\(P=|\mathcal P|\ge S^{K-1}\), and sample
\[
 p_{j,\epsilon}\qquad
 (1\le j\le m,\ \epsilon\in\{0,1\})
\]
independently and uniformly from \(\mathcal P\).  For
\(I\in\{0,1\}^m\), put \(t_I=\prod_{j=1}^m p_{j,I_j}\).
For every integer \(u\) such that
\[
 D_0\log S\le V:=\log u\le S,
 \qquad w=\min(m,V),
\]
and every integer \(1\le l\le \exp(0.005w)\), one has
\begin{equation}\label{eq:random-second-moment}
 \E\left|2^{-m}\sum_{I\in\{0,1\}^m}\e(lt_I/u)\right|^2
 \le \exp(-0.01w)
\end{equation}
for all sufficiently large \(S\).  The onset is absolute and uniform in
\(u,l,\mathcal P\).
\end{lemma}

\begin{proof}
Expanding the square gives
\[
 4^{-m}\sum_{I,J\in\{0,1\}^m}
       \E\,\e\bigl(l(t_I-t_J)/u\bigr).
\]
For uniformly counted pairs \(I,J\), their Hamming distance \(B\) has
distribution \(\operatorname{Bin}(m,1/2)\).  The exponential Markov
inequality gives
\[
 \Prob(B<m/4)
 \le 2^{m/4}\E\,2^{-B}
 =2^{m/4}(3/4)^m
 \le e^{-m/16}.
\]
We bound these pairs trivially.  Fix any remaining pair, and set
\[
 s=\left\lfloor\frac{0.75V}{K\log S}\right\rfloor.
\]
Since \(V/(K\log S)\ge1000\), and \(V\le S\), we have
\begin{equation}\label{eq:random-product-lengths}
 \frac{0.74V}{K\log S}\le s\le\frac{0.75V}{K\log S},
 \qquad 2s\le m/4
\end{equation}
for sufficiently large \(S\).

Select the first \(2s\) differing coordinates in increasing order, and
divide them into two sets of \(s\) coordinates each.  This choice depends
only on \(I,J\), not on the sampled primes.  Expose every labelled prime
sample except the \(I\)-selected sample at these \(2s\) coordinates.
Then \(t_J\) is fixed and
\[
 t_I=U_0U_1U_2,
\]
where \(U_0\) is fixed by the exposure and \(U_1,U_2\) are independent
products of \(s\) fresh uniform samples each.  Equal prime values cause
no difficulty: the \(I\)-selected and \(J\)-selected variables at a
differing coordinate are distinct independent samples.

Write
\[
 g=\gcd(lU_0,u),\qquad q=u/g.
\]
The coefficient \(lU_0/u\) reduces to a fraction with denominator
\(q\). We will show that \(q\) is usually large enough for distinct
integer values of each fresh product to remain distinct modulo \(q\).
Their small point probabilities will then give cancellation by
additive-character orthogonality.

We first bound the probability, over the exposed variables, that
\(q<u^{0.9}\).  Since \(l\le u^{0.005}\) and
\(\gcd(lU_0,u)\le l\gcd(U_0,u)\), this event implies
\(\gcd(U_0,u)>u^{0.095}\).
Call a sample contributing to \(U_0\) a hit if its prime value divides
\(u\), and count hits with their sample multiplicities.  If their number
is \(B_0\), then
\[
 \gcd(U_0,u)\le \prod_{\text{hit samples }p}p
       \le S^{(K+1)B_0}.
\]
Thus failure requires at least
\[
 k_0=\left\lceil\frac{0.08V}{(K+1)\log S}\right\rceil
\]
hits.  At most \(S\) primes in the sampling interval divide \(u\), so
each sample has hit probability at most
\(S/P\le S^{-(K-2)}\).  There are at most \(m\le S\) independent samples
in \(U_0\).  A union bound over sets of \(k_0\) positions yields
\begin{equation}\label{eq:random-gcd}
 \Prob(q<u^{0.9})
 \le S^{-(K-3)k_0}
 \le \exp\left(-\frac{0.08(K-3)}{K+1}V\right)
 \le u^{-0.05}.
\end{equation}
If \(k_0>m-2s\), the event is empty.  This argument allows prime powers
in \(u\) and repeated sampled primes.

Now fix an exposure with \(q\ge u^{0.9}\).  Every possible value of
either fresh product is at most
\[
 (2S^K)^s\le S^{(K+1)s}\le u^{0.7575}<u^{0.8}<q.
\]
By unique factorization, an integer can arise as the product of at most
\(s!\) ordered \(s\)-tuples of primes.  Consequently each atom in the
distribution of \(U_i\), for \(i=1,2\), has probability at most
\[
 \frac{s!}{P^s}\le S^{-(K-2)s}
      \le u^{-0.7252}\le u^{-0.7}.
\]
Reduction modulo \(q\) is injective on these integer values.  The
probability vectors \(\alpha,\beta\) of \(U_1,U_2\) on
\(\Z/q\Z\) therefore satisfy
\[
 \|\alpha\|_2,\|\beta\|_2\le u^{-0.35}.
\]

The reduced fraction \(lU_0/u=c/q\) has \(\gcd(c,q)=1\).
Complete additive-character orthogonality gives
\[
 \sum_{x\bmod q}\left|\sum_{y\bmod q}\beta_y\e(cxy/q)\right|^2
       =q\sum_{y\bmod q}|\beta_y|^2.
\]
Indeed the inner sum over \(x\) vanishes unless
\(q\mid c(y-y')\), equivalently \(y=y'\) modulo \(q\).
This identity holds for composite \(q\) as well.  Cauchy--Schwarz gives
\[
 \left|\sum_{x,y\bmod q}\alpha_x\beta_y\e(cxy/q)\right|
 \le \sqrt q\,\|\alpha\|_2\|\beta\|_2
 \le u^{-0.2}.
\]
The factor \(\e(-lt_J/u)\) is fixed under the exposure.  Averaging this
conditional estimate and using \eqref{eq:random-gcd} bounds the
contribution of our fixed pair \(I,J\) in absolute value by
\(u^{-0.05}+u^{-0.2}\).

Combining this with the Hamming exception, we obtain
\[
 \E\left|2^{-m}\sum_I\e(lt_I/u)\right|^2
 \le e^{-m/16}+u^{-0.05}+u^{-0.2}
 \le 3e^{-0.05w}
 \le e^{-0.01w}.
\]
The final inequality is uniform because
\(\min(m,D_0\log S)\) tends to infinity with \(S\).
\end{proof}

\subsection{A simultaneous choice of divisors}\label{subsec:selection}

We now choose all random blocks at once.  One block will suffice at all
intermediate levels outside small exceptional sets; the independent
blocks together will cover every terminal integer.

\begin{proof}[Proof of Lemma~\ref{lem:residues}]
The construction already gives the required size of \(M\), its power of
two, and divisibility of every list entry, for every realization of the
samples.  We prove the two residue properties.

Put \(H=\lfloor e^{4\eta m}\rfloor\).
At a level with \(X_j>e^S\), Lemma~\ref{lem:deterministic} and Markov's
inequality show that the number of integers
\(u\in(X_{j+1},X_j]\) for which
\[
 \left|2^{-m}\sum_{t\in T_0}\e(lCt/u)\right|>e^{-3\eta m}
\]
at some \(1\le l\le H\) is at most
\[
 X_j H e^{6\eta m}e^{-0.01m}
 \le X_j e^{-0.009m}.
\]
Conjugation changes the Fourier sign without changing the absolute
value.  Lemma~\ref{lem:discrepancy}, with scale \(m\), therefore gives at
least \(\rho|T_0|/2\) entries with
\[
 h_{j,t}(u)<\rho u\le X_{j+1}
\]
outside that exceptional set.  This holds separately at every such
level and requires no probabilistic selection.

Next suppose \(j<d\) and \(X_j\le e^S\).  Then \(X_j>e^m\), and every
\(u\in(X_{j+1},X_j]\) satisfies
\[
 (1-\eta)m<\log u\le S,\qquad
 w=\min(m,\log u)\ge(1-\eta)m.
\]
For sufficiently large \(S\), these integers satisfy the lower bound
in Lemma~\ref{lem:random}.  Moreover
\(4\eta m\le0.005w\), so that Lemma applies to all \(1\le l\le H\).
In block \(1\), Markov's inequality followed by the frequency union
bound shows that
\begin{align*}
 &\Prob\left(
   \max_{1\le l\le H}
   \left|2^{-m}\sum_{t\in T_1}\e(lt/u)\right|>e^{-3\eta m}
 \right)\\
 &\hspace{2em}\le
 \exp\bigl((10\eta-0.01(1-\eta))m\bigr)
 =e^{-0.008999m}\le e^{-0.005m}.
\end{align*}
Let \(B_j\) count integers in this level that fail this Fourier bound.
Linearity of expectation and another application of Markov give
\[
 \E B_j\le X_j e^{-0.005m},
 \qquad
 \Prob(B_j>X_j e^{-0.001m})\le e^{-0.004m}.
\]
As \(d=O(\log S)\), with probability \(1-o(1)\) these bounds hold at
every level under consideration.  No independence between different
integers or levels is needed.  Here \(C_j=1\), so
Lemma~\ref{lem:discrepancy} again gives the required
\(\rho|T_1|/2\) entries outside the exceptional sets.

Finally fix an integer \(S^{D_0}<u\le e^m\).
Now \(w=\log u\), and Lemma~\ref{lem:random}, Markov's inequality and
the union over \(1\le l\le\lfloor u^{4\eta}\rfloor\) give
\begin{align*}
 &
 \Prob\left(
  \max_{1\le l\le\lfloor u^{4\eta}\rfloor}
   \left|2^{-m}\sum_{t\in T_i}\e(lt/u)\right|>u^{-3\eta}
 \right)\\
 &\hspace{2em}\le u^{10\eta-0.01}
 =u^{-0.009}\le u^{-0.005}
\end{align*}
in each block \(i\).  On success, Lemma~\ref{lem:discrepancy}, with
scale \(\log u\), supplies a list entry for which
\[
 u\lceil t/u\rceil-t<u^{1-\eta}.
\]
The \(R=1000\) blocks are independent for this fixed \(u\).
The probability that every block fails is therefore at most \(u^{-5}\).
A union bound over all terminal integers shows that
\[
 \Prob(\text{some terminal integer has no successful block})
 \le \sum_{u>S^{D_0}}u^{-5}=o(1).
\]
Thus every terminal integer has a suitable divisor, simultaneously,
with probability tending to one.

The sum of this failure probability and the middle-level failure
probability is \(o(1)\).  Hence a realization satisfying both sets of
requirements exists.  Fix it and the resulting \(M\).
Since \(e^{-0.009m}\le e^{-0.001m}\), the same
\(c_*=0.001\) is valid at every level.  This proves both properties for
one common \(M\), for the fixed \(S\) and \(C\).
\end{proof}
